\section{An exact Gamma moment determinant}
\label{gn:sec:determinant}

The incoming ordered-boundary report asks for closed evaluations of the
positive constants $K_m$ that normalize its Gamma integral basis
\cite[Section~8 and Section~9, question~2]{repoordered}.  All of these
constants can be evaluated.  The decisive step is to remove the source's
last moment by its already proved orthogonality relation and recognize
the remaining determinant as Gustafson's symplectic Barnes integral.
We first state the parameterized determinant with its full normalization.
This is a reformulation of a classical integral, and it is credited as
such; its application to the particular source constants and the moment
identities developed below are the results of the present continuation.

For $n\ge1$ and $\mathbf a=(a_1,\ldots,a_{2n+2})$ satisfying
$\Re a_\ell>0$, put
\begin{equation}\label{gn:eq:weight}
 W_{\mathbf a}(x)=\prod_{\ell=1}^{2n+2}
       \Gamma(a_\ell+ix/2)\Gamma(a_\ell-ix/2),
 \qquad A=\sum_{\ell=1}^{2n+2}a_\ell.
\end{equation}
Define the matrix and its determinant by
\begin{equation}\label{gn:eq:matrix}
 \mathsf B_n(\mathbf a)_{r,k}
 =\int_0^\infty x^{2r-1}\sinh(k\pi x)W_{\mathbf a}(x)\dd x,
 \qquad
 \mathcal D_n(\mathbf a)=\det\mathsf B_n(\mathbf a),
 \quad 1\le r,k\le n.
\end{equation}
The sinh columns are in increasing order.  This convention is relevant
when the result is applied to a source matrix with its columns reversed.

\begin{theorem}[Parameterized Gamma moment determinant]
\label{gn:thm:determinant}
For $\Re a_\ell>0$,
\begin{equation}\label{gn:eq:determinant}
 \boxed{\displaystyle
 \mathcal D_n(\mathbf a)
 =2^{2n^2}\pi^{n(n+1)}
       \frac{\prod_{\ell<j}\Gamma(a_\ell+a_j)}{\Gamma(A)}.}
\end{equation}
Every entry is an absolutely convergent ordinary integral.  The entries,
their determinant, and every fixed mixed parameter derivative are
holomorphic on this domain and may be differentiated under the integral.
The determinant is nonzero throughout the domain and is positive when
all parameters are positive real numbers.
\end{theorem}

\begin{proof}
Stirling's formula uniformly on compact parameter sets gives
$W_{\mathbf a}(x)=O(x^C e^{-(n+1)\pi x})$ for some finite $C$.
The largest sinh factor grows as $e^{n\pi x}$.  A fixed parameter
derivative adds at most powers of $\log(2+x)$ and preserves an integrable
exponential majorant.  At zero all Gamma arguments remain in a compact
subset of the right half-plane.  These observations prove the asserted
convergence and permit the uses of Fubini and parameter differentiation.

Write
\[
 \Delta(y_1,\ldots,y_n)=\prod_{j<k}(y_k-y_j)
\]
and introduce the integral
\begin{equation}\label{gn:eq:I}
\begin{split}
 \mathcal I_n(\mathbf a)
 ={}&\int_{(0,\infty)^n}
    \Delta(x_1^2,\ldots,x_n^2)
    \Delta(\cosh\pi x_1,\ldots,\cosh\pi x_n)\\
 &\hspace{12mm}\times
       \prod_{j=1}^n x_j\sinh(\pi x_j)W_{\mathbf a}(x_j)
          \dd x_1\cdots\dd x_n.
\end{split}
\end{equation}
Expanding the two determinants in the continuous Cauchy--Binet identity
and using Fubini gives
\begin{equation}\label{gn:eq:andreief}
 \mathcal D_n(\mathbf a)
      =\frac{2^{n(n-1)/2}}{n!}\mathcal I_n(\mathbf a).
\end{equation}
Indeed, the polynomial alternant is
$\det(x_j^{2r-1})=\prod_jx_j\,\Delta(x_1^2,\ldots,x_n^2)$,
while
\[
 \det(\sinh(k\pi x_j))_{j,k=1}^n
  =2^{n(n-1)/2}\prod_j\sinh(\pi x_j)
                   \Delta(\cosh\pi x_1,\ldots,\cosh\pi x_n).
\]
The second identity follows from
$\sinh(ku)=\sinh(u)U_{k-1}(\cosh u)$ and the leading coefficient
$2^{k-1}$ of the Chebyshev polynomial $U_{k-1}$.

For precision, the classical integral used next is Gustafson's
symplectic Barnes integral \cite{gustafson1994}, in the normalization
reproduced in \cite[equation~(2)]{derkachov2017}:
\begin{equation}\label{gn:eq:gustafson}
\begin{split}
 &\int_{\mathbb R^n}
 \frac{\displaystyle\prod_{j=1}^n\prod_{\ell=1}^{2n+2}
                 \Gamma(a_\ell+iu_j)\Gamma(a_\ell-iu_j)}
 {\displaystyle\prod_j\Gamma(2iu_j)\Gamma(-2iu_j)
       \prod_{j<k}\Gamma(iu_j+iu_k)\Gamma(iu_j-iu_k)
                    \Gamma(-iu_j+iu_k)\Gamma(-iu_j-iu_k)}
                  \dd\mathbf u\\
 &\hspace{22mm}=(2\pi)^n2^nn!
       \frac{\prod_{\ell<j}\Gamma(a_\ell+a_j)}{\Gamma(A)}.
\end{split}
\end{equation}
The right-half-plane hypotheses place all numerator poles on the
appropriate sides of the imaginary integration axes in its Barnes
form.  The integrand is even in each $u_j$, so the positive-orthant
integral equals $(2\pi)^nn!$ times the Gamma quotient.

Set $x_j=2u_j$.  Reflection yields the single-variable factor
\begin{equation}\label{gn:eq:single-reflection}
 \frac1{\Gamma(ix_j)\Gamma(-ix_j)}
          =\frac{x_j\sinh(\pi x_j)}\pi,
\end{equation}
and each four-factor interaction contributes
\begin{equation}\label{gn:eq:pair-reflection}
\begin{split}
 &\left\{
 \Gamma\left(\tfrac i2(x_j+x_k)\right)
 \Gamma\left(-\tfrac i2(x_j+x_k)\right)
 \Gamma\left(\tfrac i2(x_k-x_j)\right)
 \Gamma\left(-\tfrac i2(x_k-x_j)\right)
 \right\}^{-1}\\
 &\hspace{14mm}=
 \frac{(x_k^2-x_j^2)(\cosh\pi x_k-\cosh\pi x_j)}{8\pi^2}.
\end{split}
\end{equation}
The factor $2^{-n}$ from $\dd\mathbf u$ must also be retained.
Consequently \eqref{gn:eq:gustafson} gives
\begin{equation}\label{gn:eq:I-evaluation}
 \mathcal I_n(\mathbf a)
 =2^{2n+3n(n-1)/2}\pi^{n(n+1)}n!
       \frac{\prod_{\ell<j}\Gamma(a_\ell+a_j)}{\Gamma(A)}.
\end{equation}
Combining this with \eqref{gn:eq:andreief} proves
\eqref{gn:eq:determinant}, since the two powers of $2$ add to $2n^2$.
The Gamma quotient is nonzero in the stated domain.  For positive
parameters the two Vandermonde factors in \eqref{gn:eq:I} have the same
sign, and the integrand is strictly positive away from collisions.
This proves the positivity assertion as well.
\end{proof}

The case $n=0$ may be included with the empty determinant and empty
integral both equal to one.  Its Gamma quotient is then
$\Gamma(a_1+a_2)/\Gamma(a_1+a_2)=1$.

\section{Closed evaluation of every Gamma-basis normalization}
\label{gn:sec:normalization}

We recall exactly the source data needed for the application.  For an
integer $m\ge1$, set
\begin{align}
 W_m(x)&=|\Gamma(\tfrac14+ix/2)|^{4m},
       &\phi_r(x)&=\Im(x+i/2)^{2r},\label{gn:eq:source-weight}\\
 P_j(x)&=\sinh(\pi(m-j)x)\quad(1\le j<m),
       &P_m(x)&=\tanh(\pi x),\label{gn:eq:source-multipliers}\\
 t_{r,j}&=\frac4{(2r)!}\int_0^\infty
                 \phi_r(x)P_j(x)W_m(x)\dd x,
       &T_m&=(t_{r,j})_{r,j=1}^m.\label{gn:eq:source-matrix}
\end{align}
The incoming report proves the exact top moments
\begin{equation}\label{gn:eq:source-top}
 t_{m,j}=0\quad(j<m),\qquad
 t_{m,m}=(-1)^{m+1}\frac{2^{2m+2}\pi^m}{(2m)!},
\end{equation}
by Gamma recurrence and a contour shift crossing precisely two tanh
poles.  It also proves
\begin{equation}\label{gn:eq:source-normalization}
 \det T_m=(-1)^{m(m-1)/2}K_m,\qquad K_m>0.
\end{equation}
Equations \eqref{gn:eq:source-top} and
\eqref{gn:eq:source-normalization} are inherited results, specifically
the source theorem \texttt{gb:basis}, corollary
\texttt{gb:moment-identities}, and proposition
\texttt{gb:normalization} \cite{repoordered}.  The following evaluation
answers the open question about the remaining positive constant.

\begin{theorem}[All Gamma normalization constants]
\label{gn:thm:K}
For every integer $m\ge1$,
\begin{equation}\label{gn:eq:K}
 \boxed{\displaystyle
 K_m=\frac{2^{2m^2+2}(m-1)!\,\pi^{m(4m-1)/2}}
              {\Gamma(m/2)\prod_{r=1}^m(2r)!}.}
\end{equation}
In particular,
\begin{equation}\label{gn:eq:K-values}
 K_1=8\pi,\qquad
 K_2=\frac{64}{3}\pi^7,\qquad
 K_3=\frac{16384}{135}\pi^{16}.
\end{equation}
The constants satisfy the two-step recurrence
\begin{equation}\label{gn:eq:K-recurrence}
 \boxed{\displaystyle
 \frac{K_{m+2}}{K_m}
   =\frac{2^{8m+9}(m+1)\pi^{8m+7}}
              {(2m+2)!(2m+4)!}.}
\end{equation}
Each $K_m$ is a positive rational multiple of
$\pi^{2m^2-\lceil m/2\rceil}$.
\end{theorem}

\begin{proof}
Expand $\det T_m$ along the last row, using
\eqref{gn:eq:source-top}.  Since
$\phi_r(x)=rx^{2r-1}$ plus lower odd powers, triangular row elimination
in the remaining $(m-1)\times(m-1)$ minor extracts
\[
                  \prod_{r=1}^{m-1}\frac{4r}{(2r)!}.
\]
Reversing the source sinh-column order $m-1,\ldots,1$ to the increasing
order $1,\ldots,m-1$ contributes the sign
$(-1)^{(m-1)(m-2)/2}$.  The remaining determinant is exactly
$\mathcal D_{m-1}(1/4,\ldots,1/4)$, so
\begin{equation}\label{gn:eq:detT-factorization}
\begin{split}
 \det T_m={}&(-1)^{m+1+(m-1)(m-2)/2}
 \frac{2^{2m+2}\pi^m}{(2m)!}
   \left(\prod_{r=1}^{m-1}\frac{4r}{(2r)!}\right)\\
 &\hspace{24mm}\times\mathcal D_{m-1}(1/4,\ldots,1/4).
\end{split}
\end{equation}
The sign equals $(-1)^{m(m-1)/2}$: the exponents differ by two.
Theorem~\ref{gn:thm:determinant}, with $2m$ quarter parameters, gives
\begin{equation}\label{gn:eq:quarter-minor}
 \mathcal D_{m-1}(1/4,\ldots,1/4)
 =2^{2(m-1)^2}\pi^{m(m-1)}
       \frac{\Gamma(1/2)^{m(2m-1)}}{\Gamma(m/2)}.
\end{equation}
Substitute $\Gamma(1/2)=\sqrt\pi$, compare with
\eqref{gn:eq:source-normalization}, and collect the factorials and
powers.  This yields \eqref{gn:eq:K}.  At $m=1$ the minor is empty,
so the same computation remains valid.  Taking the quotient of the
formulas at $m+2$ and $m$, and using
$\Gamma(m/2+1)=(m/2)\Gamma(m/2)$, proves
\eqref{gn:eq:K-recurrence}.  The final assertion follows from the
integer and half-integer Gamma evaluations.
\end{proof}

The result makes the inherited exact Gamma determinant representation
fully explicit.  For example, if $F_{b,m}$ is the rational fundamental
matrix of the source recurrence, its already proved identity becomes
\begin{equation}\label{gn:eq:F-explicit}
 \det F_{b,m}
 =(-1)^{m(m-1)/2}
   \frac{\Gamma(m/2)\prod_{r=1}^m(2r)!}
        {2^{2m^2+2}(m-1)!\pi^{m(4m-1)/2}}
       \det(t_{b+i,j})_{i,j=1}^m.
\end{equation}
This evaluates a previously unspecified normalization.  It preserves
the source's distinction between its eventual fixed-gap nonvanishing
theorem and its still open all-pairs reconstruction conjecture.

